\subsection{The integral class-number ratio in prime-conductor degree six}

\begin{lemma}[KT-SEXTIC-CLASS-PRODUCT]
Let $p\equiv1\pmod{12}$ be prime, let $F$ be the real cyclic sextic
subfield of $\mathbf Q(\mu_p)$, and let $F_2,F_3$ be its quadratic
and cubic subfields.  For ordinary ideal class numbers one has
\[
                         h(F_2)h(F_3)\mid h(F).
 \tag{SC.1}
\]
Thus $h(F)/(h(F_2)h(F_3))$ is a positive integer in this setting.
\end{lemma}

\begin{proof}
For $i=2,3$, let $H_i$ be the ordinary Hilbert class field of $F_i$:
it is unramified at every finite place, every real place splits, and
$\operatorname{Gal}(H_i/F_i)$ is the ordinary ideal class group.
The field $H_i$ is Galois over $\mathbf Q$, since $F_i/\mathbf Q$ is
Galois and its defining maximal unramified abelian extension is
preserved by every automorphism of $F_i$.

The unique prime over $p$ is totally ramified in every nontrivial
subextension of $F/F_i$.  Since $H_i/F_i$ is unramified there,
$H_i\cap F=F_i$.  Consequently
\[
 K_i=FH_i,\qquad [K_i:F]=h(F_i),
 \qquad \operatorname{Gal}(K_i/F)\simeq\operatorname{Gal}(H_i/F_i).
 \tag{SC.2}
\]
These are abelian extensions of $F$, unramified at all finite places
and split at its real places, and are Galois over $\mathbf Q$.

Put $E=K_2\cap K_3$ and $A=\operatorname{Gal}(E/F)$.
The extension $E/\mathbf Q$ is Galois, and $A$ is abelian.  Therefore
conjugation defines an action of
$G=\operatorname{Gal}(F/\mathbf Q)\simeq C_6$ on $A$, independent
of the choice of lifts to $\operatorname{Gal}(E/\mathbf Q)$.
This action factors through both $\operatorname{Gal}(F_2/\mathbf Q)$
and $\operatorname{Gal}(F_3/\mathbf Q)$, as follows.
If $g\in\operatorname{Gal}(F/F_i)$, the disjointness
$F\cap H_i=F_i$ gives a lift to $K_i$ acting as $g$ on $F$ and as
the identity on $H_i$.  That lift centralizes
$\operatorname{Gal}(K_i/F)$.  Its action on the quotient $A$ is
therefore trivial.  The subgroups $\operatorname{Gal}(F/F_2)$ and
$\operatorname{Gal}(F/F_3)$ have orders three and two and generate
$G$, so its action on $A$ is trivial.

Thus $A$ is central in $\Gamma=\operatorname{Gal}(E/\mathbf Q)$,
and $\Gamma/A$ is cyclic.  The group $\Gamma$ is abelian: choose
a lift $x$ of a generator of the cyclic quotient, and write every
element as $ax^j$ with $a\in A$.
The only possible finite ramified prime of $E/\mathbf Q$ is $p$.
Its inertia subgroup $I$ has order exactly six, since $E/F$ is
unramified and $F/\mathbf Q$ is totally ramified of degree six at $p$.
Because $\Gamma$ is abelian, $I$ is normal and the field $E^I$ is
unramified at every finite rational prime.  A number field with that
property has absolute discriminant one and hence is $\mathbf Q$
by Minkowski's discriminant bound.  It follows that $\Gamma=I$ and
$[E:\mathbf Q]=6$, so $E=F$.

The two Galois extensions $K_2/F$ and $K_3/F$ are consequently
linearly disjoint.  Their compositum is an abelian extension of $F$
unramified at all finite places and split at every real place, of
degree $h(F_2)h(F_3)$.  It is contained in the ordinary Hilbert
class field of $F$.  This proves (SC.1).
\end{proof}

\begin{corollary}[KT-SEXTIC-REGULATOR-RATIO-FOUR]
Use the same fields.  Let $R_2,R_3,R_6$ be their regulators, and let
$R_*$ be the primitive relative regulator in Louboutin's normalization:
for a generating relative unit with absolute logarithms
$(a,a+c,c,-a,-a-c,-c)$ it is $R_*=a^2+ac+c^2$.
Let $\chi_6$ be either primitive even Dirichlet character of order six
defining $F$, and put
\[
                       C_6=\frac{p|L(1,\chi_6)|^2}{R_*}.
 \tag{SC.3}
\]
Then $C_6$ is a positive integer divisible by four.  More precisely,
if $Q=[U_F:U_{F_2}U_{F_3}U_F^*]$, where relative units are defined
by norms lying in $\{\pm1\}$, then
\[
 Q\mid12,\qquad
 C_6=\frac{48}{Q}\frac{h(F)}{h(F_2)h(F_3)}\in4\mathbf Z_{>0}.
 \tag{SC.4}
\]
There is no hypothesis on signatures of the relative units.
\end{corollary}

\begin{proof}
Louboutin's exact regulator determinant and unit-index lemma give
$R_6=12R_2R_3R_*/Q$ with $Q\mid12$.
For clarity, the logarithmic quadratic form in that reference is
$x^2+y^2-xy$, with $x=a$ and $y=a+c$, hence equals the $R_*$ here.
The conductor-discriminant formula gives
$d_{F_2}=p$, $d_{F_3}=p^2$, and $d_F=p^5$.
The ordinary analytic class-number formula therefore reads
\[
 \begin{aligned}
 \operatorname{Res}_{s=1}\zeta_{F_2}(s)&=\frac{2h(F_2)R_2}{\sqrt p},\\
 \operatorname{Res}_{s=1}\zeta_{F_3}(s)&=\frac{4h(F_3)R_3}{p},\\
 \operatorname{Res}_{s=1}\zeta_F(s)&=\frac{32h(F)R_6}{p^{5/2}}.
 \end{aligned}
 \tag{SC.5}
\]
The Dirichlet-character factorizations give
\[
 \frac{\zeta_F(s)\zeta(s)}{\zeta_{F_2}(s)\zeta_{F_3}(s)}
       =L(s,\chi_6)L(s,\overline{\chi_6}).
\]
Taking the limit at $s=1$ and inserting (SC.5) and the regulator
determinant yields (SC.4).  Its class-number ratio is an integer
by KT-SEXTIC-CLASS-PRODUCT.  Since $Q$ divides twelve, $48/Q$ is
an integer divisible by four.  This proves the assertion.
\end{proof}

\paragraph{Exact primary locator and scope.}
S.~R.~Louboutin, \emph{Class numbers of real cyclotomic fields},
Publ. Math. Debrecen \textbf{64} (2004), printed pp.~454--455:
the relative-unit definition and regulator normalization are in
Section 3.1 on p.~454; Lemma 3 and its entire proof on p.~455
give the determinant and $Q\mid12$.
The paragraph immediately following that proof explicitly states
$h(F_2)h(F_3)\mid h(F)$ and cites Cornell--Washington, Lemma 1.
The latter paper's complete PDF was not retrieved in this audit;
(SC.1) has instead been proved above directly in the needed
prime-conductor case.  Standard Hilbert class-field existence,
base change of unramified extensions, the analytic class-number
formula, and Minkowski's bound are the stated foundational inputs.
No upper bound for $C_6$, and no exclusion of $p$ from a class group,
is asserted in this lemma and corollary alone.
